%! Author = lili
%! Date = 6/12/26

% BIS F 7 Wdh


\section{DNA-Verpackung}\label{sec:dna-verpackung}

\subsection{Histone und Nukleosomen}\label{subsec:histone-und-nukleosomen}
\defin{histon}
Fünf Klassen:
Core-Histone H2A, H2B, H3, H4 sowie Linker-Histon H1.
Unmodifizierte Histone verdichten DNA; posttranslationale Modifikationen
(Acetylierung, Methylierung, Phosphorylierung) regulieren
Chromatinzugänglichkeit und Genexpression.
\defin{nukleosom}
Benachbarte Nukleosomen sind durch Linker-DNA (10--80\,bp) verbunden und werden durch Histon H1 stabilisiert.
% F 8 - 15
\defin{linkerdna}

\frage{nukleosomverdau}

Bei verlängerter Inkubation von Chromatin oder isolierten Zellkernen mit Micrococcus-Nuklease (MNase) wird zunächst die Linker-DNA zwischen den Nukleosomen abgebaut.
Das Produkt dieser vollständigen Linker-Verdauung ist das Core-Nukleosom.
\defin{corenukleosom}
Die nukleosomale DNA selbst ist durch den Proteinkern vor weiterem Angriff geschützt, weshalb der Verdau an dieser Stelle zum Stillstand kommt.
\frage{corehistonaufbau}

Das Histon H1 bindet als Linker-Histon an der Eintrittsstelle der DNA in das Nukleosom und stabilisiert die Linker-DNA\@.
\frage{hifunktion}
Der Histonoktamer selbst ist hoch symmetrisch aufgebaut: Ein zentrales (H3₂–H4₂)-Tetramer bildet die Grundstruktur, an die beidseitig je ein H2A–H2B-Dimer bindet.
Diese Anordnung erklärt die zweifache Symmetrieachse (Pseudo-Dyade) des Nukleosoms.
\frage{histonoktameraufbau}
\defin{coreoktamer}

\begin{table}[H]
    \centering
    \begin{tabular}{ll}
        \textbf{Bezeichnung} & \textbf{Anzahl Aminosäuren} \\
        H1                   & 210--230                     \\
        H3                   & 135                         \\
        H2B                  & 125                         \\
        H2A                  & 129                         \\
        H4                   & 102
    \end{tabular}
    \caption{Protein-Bestandteile des Nukleosoms}\label{tab:table2}
\end{table}

% TODO F 16 \subsubsection*{Verschiedene Histon-Varianten}

% TODO Ist F 17 ein bsp?
Centromere sind keine Ausnahme von der Nukleosomenorganisation.
Sie haben aber ein spezialisiertes Histon, das das
kanonische H3 ersetzt: CENP-A (Centromere Protein A, auch CenH3 genannt).
CENP-A ist ein H3-Variante, die ausschließlich an Centromeren eingebaut wird.
Strukturell bildet es denselben (H3
$_2$–H4$_2$)-Tetramer-Kern wie ein normales Nukleosom, unterscheidet sich aber in der CATD-Domäne (
\textit{CENP-A targeting domain}), die für den zielgenauen Einbau verantwortlich ist.
Funktionell dient CENP-A als epigenetische Markierung des Centromers: Es definiert die Centromeridentität unabhängig
von der DNA-Sequenz.
Ein klassisches Beispiel für epigenetische Spezifizierung, denn Centromere lassen sich nicht allein über
Konsensussequenzen erklären.

CENP-A-haltige Nukleosomen rekrutieren den CCAN-Komplex (\textit{Constitutive Centromere-Associated Network}), der wiederum den Kinetochor aufbaut — die Andockstelle für Spindelmikrotubuli bei der Chromosomentrennung in der Mitose.

\frage{centromerehistone}

% TODO F 18 - 19 Architektur
% TODO F 29 - 30

\defin{zehnnmfaser}
die 10 nm-Faser besteht aus einer
kontinuierlichen Anordnung von
Nukleosomen
für Chromatin-Strukturen höherer Ordnung
ist Histon H1 nötig

% TODO abb F 21

\defin{dreißignmfaser}

% TODO abb F 22

\begin{table}[H]
    \centering
    % TODO Karteikarten
    \begin{tabular}{lll}
        & Durchmesser (nm) & bp/Windung  \\
        DNA-Doppelhelix     & 2                & 10          \\
        Nucleosomen         & 10               & 80          \\
        30 nm-Faser         & 30               & 1.200       \\
        Chromatin-Schleifen & $\sim$250        & ca.
        100.000
    \end{tabular}
    \caption{Hierarchie von DNA-Windungen}\label{tab:table3}
\end{table}

\begin{table}[H]
    \centering
    % TODO Karteikarten
    \begin{tabular}{ll}
        Nukleosomen                                & $\sim$6      \\
        30 nm-Faser                                & $\sim$40     \\
        Euchromatin                                & $\sim$1000   \\
        Heterochromatin und mitotische Chromosomen & $\sim$10 000
    \end{tabular}
    \caption{Packungsdichte von Chromatin}\label{tab:table4}
\end{table}

Chromosomen nehmen im Zellkern Chromosomenterritorien ein und sind nicht miteinander
verflochten.
Dies ist eine falschfarbige Darstellung von Chromosomenterritorien, die durch
die individuelle Färbung der Chromosomen 1–22, X und Y im Zellkern eines menschlichen Fibroblasten gewonnen wurde.
Heterochromatische Regionen, stillgelegte Gene und genarme Regionen der Chromosomen sind
typischerweise an der Peripherie des Zellkerns lokalisiert.
Aktive Gene finden sich oft an den Rändern
der Chromosomenterritorien, und aktive Gene aus mehreren Chromosomen können sich in
interchromosomalen Territorien gruppieren, die reich an Transkriptionsmaschinerie sind


\subsection{Principles of Chromatin Looping}\label{subsec:principles-of-chromatin-looping}
% TODO F 26 - 28

\subsection{Histon - Modifikationen}\label{subsec:histon---modifikationen}
% TODO F 31 - 32

\subsection{Der ``Histon - Code''}\label{subsec:der-``histon---code''}
% TODO F 33 - 34

% TODO F 35 ChIP - Chromatin immune precipitation

\subsection{Nucleosome (re)assembly}\label{subsec:nucleosome-(re)assembly}
% TODO F 40, 43

% TODO F 36 - 38  frage wie zsm hängt

% TODO F 39  frage wie zsm hängt

% TODO F 41 - 42  frage wie zsm hängt

\section{Chromatinstruktur}\label{sec:chromatinstruktur}
\defin{heterochromatin}
\defin{euchromatin}
% TODO F 47 - 48

\subsection*{Von Euchromatin zu Heterochromatin}
% TODO F 49-50

\subsection{Gene Silencing}\label{subsec:gene-silencing}
\defin{silencing}
\subsubsection{Transkriptionelles Gene Silencing (TGS)}
\defin{tgs}
(Methylierung der DNA, Heterochromatisierung)
\subsubsection{Post-transkriptionelles Gene Silencing (PTGS)}
\defin{ptgs}


% TODO F 52 - 57

\subsubsection*{RNA interference durch siRNAs}
% TODO F 58 - 62
\defin{rnai}
Wichtiges Werkzeug der reversen Genetik
zur Knockdown-Analyse von Genfunktionen.
\defin{knockdown}
z.\,B.\ via RNAi (siRNA, shRNA) oder Antisense-
Oligonukleotide.
Im Unterschied zum Knockout bleibt Restexpression
erhalten; ermöglicht Phänotypanalyse auch bei essenziellen Genen.

\subsubsection*{Antisense-Technologie}
\defin{antisensetech}

\section{microRNAs}\label{sec:micrornas}
% TODO F 63 - 69
\defin{mirna}


\subsection{Sekundärstrukturen}\label{subsec:sekundarstrukturen}
% TODO F 70

\subsection{Biogenese und Vergleich miRNAs - siRNAs}\label{subsec:biogenese-und-vergleich-mirnas---sirnas}
% TODO F 71 - 76
\defin{premirna}
\defin{primimrna}
\defin{mirnaduplex}

\subsubsection{Weitere Typen von miRNA}
% TODO F 77 - 78

\begin{bsp}
    \textbf{Überexpression von miRNAs} \\
    % TODO F 79
\end{bsp}

\subsubsection{Vorhersage von miRNA targets}
% TODO F 80

\subsubsection{Experimentelle Analyse
der miRNA-Funktionen}
% TODO F 81 - 82

\begin{bsp}
    \textbf{Überexpression von miRNA-resistenten mRNAs} \\
        % TODO F 83
\end{bsp}

\subsubsection{Endogene Rolle der miRNAs}








